% !PRINTED 497
\weberpage{497}
%% [page furniture, not body text -- source-confirmed 2026-08-10]
%% \S\ 133. \hfill Die Resolventen 7ten und 11ten Grades.
%<<CONT
so daß also die Werte von $v$ für $\omega = i$ paarweise einander
gleich werden. Aus \S\ 119, (15), (17) folgt, wenn $\omega = i$,
also $A=1$, $C=1$, $B=0$, $m=4$, $p=x=7$, $y=0$ gesetzt wird, daß
für $cc' \equiv -1 \ (\mathrm{mod}\ 7)$ $v_c = v_{c'}$ wird, so daß
also:
\weq{21}{v_\infty = v_0, \quad v_1 = v_6, \quad v_2 = v_3, \quad v_4 = v_5.}
Der Ausdruck (20) läßt sich in die beiden quadratischen Faktoren:
\[
x^2 - (1+\sqrt{7})\, x + 2, \qquad x^2 - (1-\sqrt{7})\, x + 2
\]
zerlegen, so daß die acht Werte von $x$ paarweise je einer der vier
Größen
\weq{22}{\frac{1+\sqrt{7}}{2} \pm \frac{\sqrt[4]{7}}{\sqrt{2}}, \qquad \frac{1-\sqrt{7}}{2} \pm i\, \frac{\sqrt[4]{7}}{\sqrt{2}}}
gleich werden, und es handelt sich noch darum, diese vier Werte
den einzelnen $v_\nu$ zuzuordnen. Dazu bemerken wir, daß $v_\infty$,
$v_0$ für $\omega = i$ reell sind, und daß ebenso $v_1$, $v_6$ reell
sein müssen, weil sie gleichzeitig einander gleich und konjugiert
imaginär sind.

Aus $v_0$ entsteht $v_\nu$ dadurch, daß man $q$ mit einer gewissen
Einheitswurzel multipliziert; da aber die Entwickelung von $v_0$
nach Potenzen von $q$ nur positive Koeffizienten enthält [vgl.
\S\ 24, (11)], so folgt, daß $v_0$ größer sein muß als $v_1$, und
mithin ist
\weq{23}{\begin{aligned}
\sqrt[4]{2}\, v_\infty &= \sqrt[4]{2}\, v_0 = \frac{1+\sqrt{7}}{2} + \frac{\sqrt[4]{7}}{\sqrt{2}}, \\
\sqrt[4]{2}\, v_1 &= \sqrt[4]{2}\, v_6 = \frac{1+\sqrt{7}}{2} - \frac{\sqrt[4]{7}}{\sqrt{2}}.
\end{aligned}}
Für die beiden anderen Wurzelpaare ergibt sich:
\weq{24}{\begin{aligned}
\sqrt[4]{2}\, v_2 &= \sqrt{2}\, v_3 = \frac{1-\sqrt{7}}{2} + i\, \frac{\sqrt[4]{7}}{\sqrt{2}}, \\ % [sic] radical on $v_3$ printed as $\sqrt{2}$, not $\sqrt[4]{2}$
\sqrt[4]{2}\, v_4 &= \sqrt{2}\, v_5 = \frac{1-\sqrt{7}}{2} - i\, \frac{\sqrt[4]{7}}{\sqrt{2}}. % [sic] radical on $v_5$ printed as $\sqrt{2}$, not $\sqrt[4]{2}$
\end{aligned}}
Daß in diesen Ausdrücken das Vorzeichen von $i$ richtig gewählt
ist, schließt man auf folgende Weise:

Aus (23) ist zu ersehen, daß für $\omega = i$, also für
$u = \sqrt{2}$, die Wurzel $w_0$ verschwindet, und daß also für
diesen Wert $A_7$ verschwinden muß. Nach den oben nachgewiesenen
Eigenschaften ist hierdurch $A_7$ vollständig bestimmt, nämlich:
% [sic] printed $u=\sqrt2$ conflicts with $u=f(i)=\sqrt[4]{2}$ on the preceding page
\weq{25}{A_7 = -\left( u^{12} - \frac{64}{u^{12}} \right)^2,}
%CONT>>
